% codex_solutions.tex — 工作區共用的習題詳解版型（2026-10-09 從量導 v3 詳解抽出）
% 用於：量導 quantum_exercises_v3.sty、Neamen neamen_exercises_v4.sty、熱物理 thermal_solutions_v2.sty。
% 各專案的 .sty 只放記號巨集與版面差異，並用下面這行載入本檔：
%   \IfFileExists{codex_solutions.tex}{% codex_solutions.tex — 工作區共用的習題詳解版型（2026-10-09 從量導 v3 詳解抽出）
% 用於：量導 quantum_exercises_v3.sty、Neamen neamen_exercises_v4.sty、熱物理 thermal_solutions_v2.sty。
% 各專案的 .sty 只放記號巨集與版面差異，並用下面這行載入本檔：
%   \IfFileExists{codex_solutions.tex}{% codex_solutions.tex — 工作區共用的習題詳解版型（2026-10-09 從量導 v3 詳解抽出）
% 用於：量導 quantum_exercises_v3.sty、Neamen neamen_exercises_v4.sty、熱物理 thermal_solutions_v2.sty。
% 各專案的 .sty 只放記號巨集與版面差異，並用下面這行載入本檔：
%   \IfFileExists{codex_solutions.tex}{% codex_solutions.tex — 工作區共用的習題詳解版型（2026-10-09 從量導 v3 詳解抽出）
% 用於：量導 quantum_exercises_v3.sty、Neamen neamen_exercises_v4.sty、熱物理 thermal_solutions_v2.sty。
% 各專案的 .sty 只放記號巨集與版面差異，並用下面這行載入本檔：
%   \IfFileExists{codex_solutions.tex}{\input{codex_solutions.tex}}{\input{"D:/Codex workplace/_templates/styles/codex_solutions.tex"}}
% （Windows 走後者的絕對路徑；Linux 沙箱的 texbuild.sh 與雲端會把本資料夾放進 TEXMFHOME，走前者。）
% 版面：A4、10pt（主檔 \documentclass[10pt,a4paper]{article}）、四邊 2cm、行距 1.3、Noto Sans TC；
% 深藍標題、淺灰英文題目框（TeX Gyre Termes）、淺橘結果框。字型設定在建置資料夾的 fonts.tex。
\RequirePackage[no-math]{fontspec}
\RequirePackage{xeCJK}
% 英文撇號改回西文字元並取比例寬字形（原因見 codex_notes.tex 同一段；英文題目框的 Termes 本來就是比例寬）。2026-10-22
\defaultfontfeatures+[\rmfamily]{RawFeature=+pwid}
\xeCJKDeclareCharClass{Default}{"2018,"2019}
\input{fonts.tex}
\RequirePackage{amsmath,amssymb,amsthm,mathtools}
\RequirePackage[a4paper,margin=2cm]{geometry}
\RequirePackage{graphicx,float,booktabs,longtable,array,tabularx}  % float 要在 setspace 之前載入（[H] 圖內單倍行距）
\RequirePackage{enumitem,setspace,xcolor,fancyhdr,needspace}
\RequirePackage[most]{tcolorbox}
\RequirePackage{titlesec,titletoc}
% 專案要在 hyperref 之前多載的套件：載入本檔前 \def\codexsolutionsprehyperref{\RequirePackage{...}}
\csname codexsolutionsprehyperref\endcsname
\RequirePackage[unicode]{hyperref}

% ---- 色票 ----
\definecolor{qblue}{HTML}{1D5FA8}    % 章標題、目錄、條列符號
\definecolor{qnavy}{HTML}{173D6B}    % 題目標題、小標
\definecolor{qorange}{HTML}{EB6834}  % 結果框
\definecolor{qorangebg}{HTML}{FFF4E8}
\definecolor{qgraybg}{HTML}{F3F4F6}  % 原文框
\definecolor{qgraybar}{HTML}{8A929C}
\definecolor{qrule}{HTML}{D8DDE3}    % 細分隔線
\definecolor{qmuted}{HTML}{6B7280}   % 次要文字
\definecolor{qtext}{HTML}{1F2329}    % 內文
\hypersetup{colorlinks=true,linkcolor=qnavy,urlcolor=qblue,citecolor=qblue}

% ---- 字型：英文題目用襯線字，與中文詳解區隔 ----
\newfontfamily\statementfont{texgyretermes-regular.otf}[
  BoldFont=texgyretermes-bold.otf,ItalicFont=texgyretermes-italic.otf,
  BoldItalicFont=texgyretermes-bolditalic.otf,Scale=1.04,Ligatures=TeX]

% ---- 段落 ----
\setstretch{1.3}
\setlength{\parindent}{0pt}
\setlength{\parskip}{0.55em}
\setlength{\headheight}{16pt}
\setlength{\headsep}{0.45cm}
\setlength{\footskip}{1cm}
\setlist{itemsep=0.2em,topsep=0.3em}
\setlist[itemize]{label=\textcolor{qblue}{\small\textbullet}}
\setlist[enumerate]{label=\textcolor{qblue}{\arabic*.}}
\renewcommand{\arraystretch}{1.16}
\allowdisplaybreaks[2]
\raggedbottom
\emergencystretch=3em
\AtBeginDocument{\color{qtext}}

% ---- 頁首頁尾 ----
\pagestyle{fancy}
\fancyhf{}
\fancyfoot[C]{\color{qmuted}\small\thepage}
\renewcommand{\headrulewidth}{0pt}
\newcommand{\courseheader}[2]{%
  \fancyhead[L]{\color{qmuted}\small #1}\fancyhead[R]{\color{qmuted}\small #2}}

% ---- 標題 ----
% 題目標題：\problem{2.8}{指數波函數的動量空間表示}
\titleformat{\section}[block]{\color{qnavy}\large\bfseries\boldmath}{}{0pt}{}
\titlespacing*{\section}{0pt}{2.4em}{0.7em}
% 詳解內的小標（思路、步驟一、檢查……）
\titleformat{\subsection}{\color{qnavy}\bfseries\boldmath}{}{0pt}{}
\titlespacing*{\subsection}{0pt}{1.1em}{0.2em}
\titleformat{\subsubsection}{\color{qnavy}\bfseries\boldmath}{}{0pt}{}
\titlespacing*{\subsubsection}{0pt}{0.8em}{0.2em}
\titleformat{\paragraph}[runin]{\color{qnavy}\bfseries}{}{0pt}{}
\titlespacing*{\paragraph}{0pt}{0.6em}{0.6em}

\newcommand{\problem}[2]{%
  \Needspace{10\baselineskip}%
  \section*{{\color{qmuted}\mdseries Problem #1}\hspace{0.8em}#2}%
  \addcontentsline{toc}{subsection}{Problem #1\hspace{0.8em}#2}}

% 章首：\chapterhead{CHAPTER 2}{第 2 章　波函數與機率}
\newcommand{\chapterhead}[2]{%
  \clearpage
  {\color{qmuted}\small #1\par}\vspace{2pt}
  {\color{qnavy}\fontsize{16pt}{21pt}\selectfont\bfseries #2\par}\vspace{6pt}
  {\color{qrule}\rule{\linewidth}{0.6pt}\par}\vspace{0.4em}
  \addcontentsline{toc}{section}{#2}}

% ---- 目錄 ----
\renewcommand{\contentsname}{目錄}
\titlecontents{section}[0pt]{\addvspace{0.8em}\color{qblue}\bfseries}{}{}
  {\titlerule*[0.6pc]{.}\contentspage}
\titlecontents{subsection}[1.5em]{\color{qtext}}{}{}
  {\color{qmuted}\titlerule*[0.6pc]{.}\contentspage}
% 頁碼欄放得下三位數（article 預設 1.55em，詳解冊過 100 頁時每行目錄溢出 1–2pt；DE CH10、CH11）。2026-10-27
\makeatletter
\renewcommand{\@pnumwidth}{2.2em}
\renewcommand{\@tocrmarg}{3em}
\makeatother

% ---- 框 ----
\newenvironment{chapterintro}{\par\color{qmuted}}{\par\vspace{0.4em}}

% 題目（英文原題）。跨頁時續段右上角標「（續）」，免得看成新的一題（DE CH11 版面審查）。2026-10-27
\newtcolorbox{statement}{enhanced,breakable,frame hidden,arc=2pt,
  colback=qgraybg,left=10pt,right=10pt,top=6pt,bottom=6pt,
  before skip=0.2em,after skip=0.9em,before upper={\statementfont},
  overlay middle and last={\node[anchor=north east,inner sep=0pt,font=\scriptsize,text=qmuted]
    at ([xshift=-6pt,yshift=-4pt]frame.north east) {（續）};}}

% 結果框；預設無標籤，需要時 \begin{resultbox}[答案 (b)]
\newtcolorbox{resultbox}[1][]{enhanced,frame hidden,sharp corners,
  colback=qorangebg,borderline west={2pt}{0pt}{qorange},
  left=10pt,right=10pt,top=5pt,bottom=5pt,
  before skip=0.6em,after skip=0.8em,
  before upper={\if\relax\detokenize{#1}\relax\else{\small\bfseries\color{qorange}#1}\par\fi}}

% 提醒或補充
\newtcolorbox{notebox}[1][說明]{enhanced,breakable,frame hidden,sharp corners,
  colback=qblue!5,borderline west={2pt}{0pt}{qblue},
  left=11pt,right=11pt,top=6pt,bottom=6pt,
  before upper={{\small\bfseries\color{qblue}#1}\par}}

\tcbset{highlight math style={enhanced,frame hidden,sharp corners,colback=qorangebg,
  borderline west={2pt}{0pt}{qorange},left=6pt,right=6pt,top=3pt,bottom=3pt}}

% ---- 數學巨集（各專案共用的兩個；其餘記號在專案的 .sty）----
\newcommand{\dd}{\mathop{}\!\mathrm{d}}
\newcommand{\e}{\mathrm{e}}
% \keyresult{...} 的內容當文字排（數學用 $...$），放進結果框；已在數學環境內時改為橘色底
\newcommand{\keyresult}[1]{%
  \ifmmode\tcbhighmath{\displaystyle #1}%
  \else\begin{resultbox}{\centering #1\par}\end{resultbox}\fi}
}{\input{"D:/Codex workplace/_templates/styles/codex_solutions.tex"}}
% （Windows 走後者的絕對路徑；Linux 沙箱的 texbuild.sh 與雲端會把本資料夾放進 TEXMFHOME，走前者。）
% 版面：A4、10pt（主檔 \documentclass[10pt,a4paper]{article}）、四邊 2cm、行距 1.3、Noto Sans TC；
% 深藍標題、淺灰英文題目框（TeX Gyre Termes）、淺橘結果框。字型設定在建置資料夾的 fonts.tex。
\RequirePackage[no-math]{fontspec}
\RequirePackage{xeCJK}
% 英文撇號改回西文字元並取比例寬字形（原因見 codex_notes.tex 同一段；英文題目框的 Termes 本來就是比例寬）。2026-10-22
\defaultfontfeatures+[\rmfamily]{RawFeature=+pwid}
\xeCJKDeclareCharClass{Default}{"2018,"2019}
\input{fonts.tex}
\RequirePackage{amsmath,amssymb,amsthm,mathtools}
\RequirePackage[a4paper,margin=2cm]{geometry}
\RequirePackage{graphicx,float,booktabs,longtable,array,tabularx}  % float 要在 setspace 之前載入（[H] 圖內單倍行距）
\RequirePackage{enumitem,setspace,xcolor,fancyhdr,needspace}
\RequirePackage[most]{tcolorbox}
\RequirePackage{titlesec,titletoc}
% 專案要在 hyperref 之前多載的套件：載入本檔前 \def\codexsolutionsprehyperref{\RequirePackage{...}}
\csname codexsolutionsprehyperref\endcsname
\RequirePackage[unicode]{hyperref}

% ---- 色票 ----
\definecolor{qblue}{HTML}{1D5FA8}    % 章標題、目錄、條列符號
\definecolor{qnavy}{HTML}{173D6B}    % 題目標題、小標
\definecolor{qorange}{HTML}{EB6834}  % 結果框
\definecolor{qorangebg}{HTML}{FFF4E8}
\definecolor{qgraybg}{HTML}{F3F4F6}  % 原文框
\definecolor{qgraybar}{HTML}{8A929C}
\definecolor{qrule}{HTML}{D8DDE3}    % 細分隔線
\definecolor{qmuted}{HTML}{6B7280}   % 次要文字
\definecolor{qtext}{HTML}{1F2329}    % 內文
\hypersetup{colorlinks=true,linkcolor=qnavy,urlcolor=qblue,citecolor=qblue}

% ---- 字型：英文題目用襯線字，與中文詳解區隔 ----
\newfontfamily\statementfont{texgyretermes-regular.otf}[
  BoldFont=texgyretermes-bold.otf,ItalicFont=texgyretermes-italic.otf,
  BoldItalicFont=texgyretermes-bolditalic.otf,Scale=1.04,Ligatures=TeX]

% ---- 段落 ----
\setstretch{1.3}
\setlength{\parindent}{0pt}
\setlength{\parskip}{0.55em}
\setlength{\headheight}{16pt}
\setlength{\headsep}{0.45cm}
\setlength{\footskip}{1cm}
\setlist{itemsep=0.2em,topsep=0.3em}
\setlist[itemize]{label=\textcolor{qblue}{\small\textbullet}}
\setlist[enumerate]{label=\textcolor{qblue}{\arabic*.}}
\renewcommand{\arraystretch}{1.16}
\allowdisplaybreaks[2]
\raggedbottom
\emergencystretch=3em
\AtBeginDocument{\color{qtext}}

% ---- 頁首頁尾 ----
\pagestyle{fancy}
\fancyhf{}
\fancyfoot[C]{\color{qmuted}\small\thepage}
\renewcommand{\headrulewidth}{0pt}
\newcommand{\courseheader}[2]{%
  \fancyhead[L]{\color{qmuted}\small #1}\fancyhead[R]{\color{qmuted}\small #2}}

% ---- 標題 ----
% 題目標題：\problem{2.8}{指數波函數的動量空間表示}
\titleformat{\section}[block]{\color{qnavy}\large\bfseries\boldmath}{}{0pt}{}
\titlespacing*{\section}{0pt}{2.4em}{0.7em}
% 詳解內的小標（思路、步驟一、檢查……）
\titleformat{\subsection}{\color{qnavy}\bfseries\boldmath}{}{0pt}{}
\titlespacing*{\subsection}{0pt}{1.1em}{0.2em}
\titleformat{\subsubsection}{\color{qnavy}\bfseries\boldmath}{}{0pt}{}
\titlespacing*{\subsubsection}{0pt}{0.8em}{0.2em}
\titleformat{\paragraph}[runin]{\color{qnavy}\bfseries}{}{0pt}{}
\titlespacing*{\paragraph}{0pt}{0.6em}{0.6em}

\newcommand{\problem}[2]{%
  \Needspace{10\baselineskip}%
  \section*{{\color{qmuted}\mdseries Problem #1}\hspace{0.8em}#2}%
  \addcontentsline{toc}{subsection}{Problem #1\hspace{0.8em}#2}}

% 章首：\chapterhead{CHAPTER 2}{第 2 章　波函數與機率}
\newcommand{\chapterhead}[2]{%
  \clearpage
  {\color{qmuted}\small #1\par}\vspace{2pt}
  {\color{qnavy}\fontsize{16pt}{21pt}\selectfont\bfseries #2\par}\vspace{6pt}
  {\color{qrule}\rule{\linewidth}{0.6pt}\par}\vspace{0.4em}
  \addcontentsline{toc}{section}{#2}}

% ---- 目錄 ----
\renewcommand{\contentsname}{目錄}
\titlecontents{section}[0pt]{\addvspace{0.8em}\color{qblue}\bfseries}{}{}
  {\titlerule*[0.6pc]{.}\contentspage}
\titlecontents{subsection}[1.5em]{\color{qtext}}{}{}
  {\color{qmuted}\titlerule*[0.6pc]{.}\contentspage}
% 頁碼欄放得下三位數（article 預設 1.55em，詳解冊過 100 頁時每行目錄溢出 1–2pt；DE CH10、CH11）。2026-10-27
\makeatletter
\renewcommand{\@pnumwidth}{2.2em}
\renewcommand{\@tocrmarg}{3em}
\makeatother

% ---- 框 ----
\newenvironment{chapterintro}{\par\color{qmuted}}{\par\vspace{0.4em}}

% 題目（英文原題）。跨頁時續段右上角標「（續）」，免得看成新的一題（DE CH11 版面審查）。2026-10-27
\newtcolorbox{statement}{enhanced,breakable,frame hidden,arc=2pt,
  colback=qgraybg,left=10pt,right=10pt,top=6pt,bottom=6pt,
  before skip=0.2em,after skip=0.9em,before upper={\statementfont},
  overlay middle and last={\node[anchor=north east,inner sep=0pt,font=\scriptsize,text=qmuted]
    at ([xshift=-6pt,yshift=-4pt]frame.north east) {（續）};}}

% 結果框；預設無標籤，需要時 \begin{resultbox}[答案 (b)]
\newtcolorbox{resultbox}[1][]{enhanced,frame hidden,sharp corners,
  colback=qorangebg,borderline west={2pt}{0pt}{qorange},
  left=10pt,right=10pt,top=5pt,bottom=5pt,
  before skip=0.6em,after skip=0.8em,
  before upper={\if\relax\detokenize{#1}\relax\else{\small\bfseries\color{qorange}#1}\par\fi}}

% 提醒或補充
\newtcolorbox{notebox}[1][說明]{enhanced,breakable,frame hidden,sharp corners,
  colback=qblue!5,borderline west={2pt}{0pt}{qblue},
  left=11pt,right=11pt,top=6pt,bottom=6pt,
  before upper={{\small\bfseries\color{qblue}#1}\par}}

\tcbset{highlight math style={enhanced,frame hidden,sharp corners,colback=qorangebg,
  borderline west={2pt}{0pt}{qorange},left=6pt,right=6pt,top=3pt,bottom=3pt}}

% ---- 數學巨集（各專案共用的兩個；其餘記號在專案的 .sty）----
\newcommand{\dd}{\mathop{}\!\mathrm{d}}
\newcommand{\e}{\mathrm{e}}
% \keyresult{...} 的內容當文字排（數學用 $...$），放進結果框；已在數學環境內時改為橘色底
\newcommand{\keyresult}[1]{%
  \ifmmode\tcbhighmath{\displaystyle #1}%
  \else\begin{resultbox}{\centering #1\par}\end{resultbox}\fi}
}{\input{"D:/Codex workplace/_templates/styles/codex_solutions.tex"}}
% （Windows 走後者的絕對路徑；Linux 沙箱的 texbuild.sh 與雲端會把本資料夾放進 TEXMFHOME，走前者。）
% 版面：A4、10pt（主檔 \documentclass[10pt,a4paper]{article}）、四邊 2cm、行距 1.3、Noto Sans TC；
% 深藍標題、淺灰英文題目框（TeX Gyre Termes）、淺橘結果框。字型設定在建置資料夾的 fonts.tex。
\RequirePackage[no-math]{fontspec}
\RequirePackage{xeCJK}
% 英文撇號改回西文字元並取比例寬字形（原因見 codex_notes.tex 同一段；英文題目框的 Termes 本來就是比例寬）。2026-10-22
\defaultfontfeatures+[\rmfamily]{RawFeature=+pwid}
\xeCJKDeclareCharClass{Default}{"2018,"2019}
\input{fonts.tex}
\RequirePackage{amsmath,amssymb,amsthm,mathtools}
\RequirePackage[a4paper,margin=2cm]{geometry}
\RequirePackage{graphicx,float,booktabs,longtable,array,tabularx}  % float 要在 setspace 之前載入（[H] 圖內單倍行距）
\RequirePackage{enumitem,setspace,xcolor,fancyhdr,needspace}
\RequirePackage[most]{tcolorbox}
\RequirePackage{titlesec,titletoc}
% 專案要在 hyperref 之前多載的套件：載入本檔前 \def\codexsolutionsprehyperref{\RequirePackage{...}}
\csname codexsolutionsprehyperref\endcsname
\RequirePackage[unicode]{hyperref}

% ---- 色票 ----
\definecolor{qblue}{HTML}{1D5FA8}    % 章標題、目錄、條列符號
\definecolor{qnavy}{HTML}{173D6B}    % 題目標題、小標
\definecolor{qorange}{HTML}{EB6834}  % 結果框
\definecolor{qorangebg}{HTML}{FFF4E8}
\definecolor{qgraybg}{HTML}{F3F4F6}  % 原文框
\definecolor{qgraybar}{HTML}{8A929C}
\definecolor{qrule}{HTML}{D8DDE3}    % 細分隔線
\definecolor{qmuted}{HTML}{6B7280}   % 次要文字
\definecolor{qtext}{HTML}{1F2329}    % 內文
\hypersetup{colorlinks=true,linkcolor=qnavy,urlcolor=qblue,citecolor=qblue}

% ---- 字型：英文題目用襯線字，與中文詳解區隔 ----
\newfontfamily\statementfont{texgyretermes-regular.otf}[
  BoldFont=texgyretermes-bold.otf,ItalicFont=texgyretermes-italic.otf,
  BoldItalicFont=texgyretermes-bolditalic.otf,Scale=1.04,Ligatures=TeX]

% ---- 段落 ----
\setstretch{1.3}
\setlength{\parindent}{0pt}
\setlength{\parskip}{0.55em}
\setlength{\headheight}{16pt}
\setlength{\headsep}{0.45cm}
\setlength{\footskip}{1cm}
\setlist{itemsep=0.2em,topsep=0.3em}
\setlist[itemize]{label=\textcolor{qblue}{\small\textbullet}}
\setlist[enumerate]{label=\textcolor{qblue}{\arabic*.}}
\renewcommand{\arraystretch}{1.16}
\allowdisplaybreaks[2]
\raggedbottom
\emergencystretch=3em
\AtBeginDocument{\color{qtext}}

% ---- 頁首頁尾 ----
\pagestyle{fancy}
\fancyhf{}
\fancyfoot[C]{\color{qmuted}\small\thepage}
\renewcommand{\headrulewidth}{0pt}
\newcommand{\courseheader}[2]{%
  \fancyhead[L]{\color{qmuted}\small #1}\fancyhead[R]{\color{qmuted}\small #2}}

% ---- 標題 ----
% 題目標題：\problem{2.8}{指數波函數的動量空間表示}
\titleformat{\section}[block]{\color{qnavy}\large\bfseries\boldmath}{}{0pt}{}
\titlespacing*{\section}{0pt}{2.4em}{0.7em}
% 詳解內的小標（思路、步驟一、檢查……）
\titleformat{\subsection}{\color{qnavy}\bfseries\boldmath}{}{0pt}{}
\titlespacing*{\subsection}{0pt}{1.1em}{0.2em}
\titleformat{\subsubsection}{\color{qnavy}\bfseries\boldmath}{}{0pt}{}
\titlespacing*{\subsubsection}{0pt}{0.8em}{0.2em}
\titleformat{\paragraph}[runin]{\color{qnavy}\bfseries}{}{0pt}{}
\titlespacing*{\paragraph}{0pt}{0.6em}{0.6em}

\newcommand{\problem}[2]{%
  \Needspace{10\baselineskip}%
  \section*{{\color{qmuted}\mdseries Problem #1}\hspace{0.8em}#2}%
  \addcontentsline{toc}{subsection}{Problem #1\hspace{0.8em}#2}}

% 章首：\chapterhead{CHAPTER 2}{第 2 章　波函數與機率}
\newcommand{\chapterhead}[2]{%
  \clearpage
  {\color{qmuted}\small #1\par}\vspace{2pt}
  {\color{qnavy}\fontsize{16pt}{21pt}\selectfont\bfseries #2\par}\vspace{6pt}
  {\color{qrule}\rule{\linewidth}{0.6pt}\par}\vspace{0.4em}
  \addcontentsline{toc}{section}{#2}}

% ---- 目錄 ----
\renewcommand{\contentsname}{目錄}
\titlecontents{section}[0pt]{\addvspace{0.8em}\color{qblue}\bfseries}{}{}
  {\titlerule*[0.6pc]{.}\contentspage}
\titlecontents{subsection}[1.5em]{\color{qtext}}{}{}
  {\color{qmuted}\titlerule*[0.6pc]{.}\contentspage}
% 頁碼欄放得下三位數（article 預設 1.55em，詳解冊過 100 頁時每行目錄溢出 1–2pt；DE CH10、CH11）。2026-10-27
\makeatletter
\renewcommand{\@pnumwidth}{2.2em}
\renewcommand{\@tocrmarg}{3em}
\makeatother

% ---- 框 ----
\newenvironment{chapterintro}{\par\color{qmuted}}{\par\vspace{0.4em}}

% 題目（英文原題）。跨頁時續段右上角標「（續）」，免得看成新的一題（DE CH11 版面審查）。2026-10-27
\newtcolorbox{statement}{enhanced,breakable,frame hidden,arc=2pt,
  colback=qgraybg,left=10pt,right=10pt,top=6pt,bottom=6pt,
  before skip=0.2em,after skip=0.9em,before upper={\statementfont},
  overlay middle and last={\node[anchor=north east,inner sep=0pt,font=\scriptsize,text=qmuted]
    at ([xshift=-6pt,yshift=-4pt]frame.north east) {（續）};}}

% 結果框；預設無標籤，需要時 \begin{resultbox}[答案 (b)]
\newtcolorbox{resultbox}[1][]{enhanced,frame hidden,sharp corners,
  colback=qorangebg,borderline west={2pt}{0pt}{qorange},
  left=10pt,right=10pt,top=5pt,bottom=5pt,
  before skip=0.6em,after skip=0.8em,
  before upper={\if\relax\detokenize{#1}\relax\else{\small\bfseries\color{qorange}#1}\par\fi}}

% 提醒或補充
\newtcolorbox{notebox}[1][說明]{enhanced,breakable,frame hidden,sharp corners,
  colback=qblue!5,borderline west={2pt}{0pt}{qblue},
  left=11pt,right=11pt,top=6pt,bottom=6pt,
  before upper={{\small\bfseries\color{qblue}#1}\par}}

\tcbset{highlight math style={enhanced,frame hidden,sharp corners,colback=qorangebg,
  borderline west={2pt}{0pt}{qorange},left=6pt,right=6pt,top=3pt,bottom=3pt}}

% ---- 數學巨集（各專案共用的兩個；其餘記號在專案的 .sty）----
\newcommand{\dd}{\mathop{}\!\mathrm{d}}
\newcommand{\e}{\mathrm{e}}
% \keyresult{...} 的內容當文字排（數學用 $...$），放進結果框；已在數學環境內時改為橘色底
\newcommand{\keyresult}[1]{%
  \ifmmode\tcbhighmath{\displaystyle #1}%
  \else\begin{resultbox}{\centering #1\par}\end{resultbox}\fi}
}{\input{"D:/Codex workplace/_templates/styles/codex_solutions.tex"}}
% （Windows 走後者的絕對路徑；Linux 沙箱的 texbuild.sh 與雲端會把本資料夾放進 TEXMFHOME，走前者。）
% 版面：A4、10pt（主檔 \documentclass[10pt,a4paper]{article}）、四邊 2cm、行距 1.3、Noto Sans TC；
% 深藍標題、淺灰英文題目框（TeX Gyre Termes）、淺橘結果框。字型設定在建置資料夾的 fonts.tex。
\RequirePackage[no-math]{fontspec}
\RequirePackage{xeCJK}
% 英文撇號改回西文字元並取比例寬字形（原因見 codex_notes.tex 同一段；英文題目框的 Termes 本來就是比例寬）。2026-10-22
\defaultfontfeatures+[\rmfamily]{RawFeature=+pwid}
\xeCJKDeclareCharClass{Default}{"2018,"2019}
\input{fonts.tex}
\RequirePackage{amsmath,amssymb,amsthm,mathtools}
\RequirePackage[a4paper,margin=2cm]{geometry}
\RequirePackage{graphicx,float,booktabs,longtable,array,tabularx}  % float 要在 setspace 之前載入（[H] 圖內單倍行距）
\RequirePackage{enumitem,setspace,xcolor,fancyhdr,needspace}
\RequirePackage[most]{tcolorbox}
\RequirePackage{titlesec,titletoc}
% 專案要在 hyperref 之前多載的套件：載入本檔前 \def\codexsolutionsprehyperref{\RequirePackage{...}}
\csname codexsolutionsprehyperref\endcsname
\RequirePackage[unicode]{hyperref}

% ---- 色票 ----
\definecolor{qblue}{HTML}{1D5FA8}    % 章標題、目錄、條列符號
\definecolor{qnavy}{HTML}{173D6B}    % 題目標題、小標
\definecolor{qorange}{HTML}{EB6834}  % 結果框
\definecolor{qorangebg}{HTML}{FFF4E8}
\definecolor{qgraybg}{HTML}{F3F4F6}  % 原文框
\definecolor{qgraybar}{HTML}{8A929C}
\definecolor{qrule}{HTML}{D8DDE3}    % 細分隔線
\definecolor{qmuted}{HTML}{6B7280}   % 次要文字
\definecolor{qtext}{HTML}{1F2329}    % 內文
\hypersetup{colorlinks=true,linkcolor=qnavy,urlcolor=qblue,citecolor=qblue}

% ---- 字型：英文題目用襯線字，與中文詳解區隔 ----
\newfontfamily\statementfont{texgyretermes-regular.otf}[
  BoldFont=texgyretermes-bold.otf,ItalicFont=texgyretermes-italic.otf,
  BoldItalicFont=texgyretermes-bolditalic.otf,Scale=1.04,Ligatures=TeX]

% ---- 段落 ----
\setstretch{1.3}
\setlength{\parindent}{0pt}
\setlength{\parskip}{0.55em}
\setlength{\headheight}{16pt}
\setlength{\headsep}{0.45cm}
\setlength{\footskip}{1cm}
\setlist{itemsep=0.2em,topsep=0.3em}
\setlist[itemize]{label=\textcolor{qblue}{\small\textbullet}}
\setlist[enumerate]{label=\textcolor{qblue}{\arabic*.}}
\renewcommand{\arraystretch}{1.16}
\allowdisplaybreaks[2]
\raggedbottom
\emergencystretch=3em
\AtBeginDocument{\color{qtext}}

% ---- 頁首頁尾 ----
\pagestyle{fancy}
\fancyhf{}
\fancyfoot[C]{\color{qmuted}\small\thepage}
\renewcommand{\headrulewidth}{0pt}
\newcommand{\courseheader}[2]{%
  \fancyhead[L]{\color{qmuted}\small #1}\fancyhead[R]{\color{qmuted}\small #2}}

% ---- 標題 ----
% 題目標題：\problem{2.8}{指數波函數的動量空間表示}
\titleformat{\section}[block]{\color{qnavy}\large\bfseries\boldmath}{}{0pt}{}
\titlespacing*{\section}{0pt}{2.4em}{0.7em}
% 詳解內的小標（思路、步驟一、檢查……）
\titleformat{\subsection}{\color{qnavy}\bfseries\boldmath}{}{0pt}{}
\titlespacing*{\subsection}{0pt}{1.1em}{0.2em}
\titleformat{\subsubsection}{\color{qnavy}\bfseries\boldmath}{}{0pt}{}
\titlespacing*{\subsubsection}{0pt}{0.8em}{0.2em}
\titleformat{\paragraph}[runin]{\color{qnavy}\bfseries}{}{0pt}{}
\titlespacing*{\paragraph}{0pt}{0.6em}{0.6em}

\newcommand{\problem}[2]{%
  \Needspace{10\baselineskip}%
  \section*{{\color{qmuted}\mdseries Problem #1}\hspace{0.8em}#2}%
  \addcontentsline{toc}{subsection}{Problem #1\hspace{0.8em}#2}}

% 章首：\chapterhead{CHAPTER 2}{第 2 章　波函數與機率}
\newcommand{\chapterhead}[2]{%
  \clearpage
  {\color{qmuted}\small #1\par}\vspace{2pt}
  {\color{qnavy}\fontsize{16pt}{21pt}\selectfont\bfseries #2\par}\vspace{6pt}
  {\color{qrule}\rule{\linewidth}{0.6pt}\par}\vspace{0.4em}
  \addcontentsline{toc}{section}{#2}}

% ---- 目錄 ----
\renewcommand{\contentsname}{目錄}
\titlecontents{section}[0pt]{\addvspace{0.8em}\color{qblue}\bfseries}{}{}
  {\titlerule*[0.6pc]{.}\contentspage}
\titlecontents{subsection}[1.5em]{\color{qtext}}{}{}
  {\color{qmuted}\titlerule*[0.6pc]{.}\contentspage}
% 頁碼欄放得下三位數（article 預設 1.55em，詳解冊過 100 頁時每行目錄溢出 1–2pt；DE CH10、CH11）。2026-10-27
\makeatletter
\renewcommand{\@pnumwidth}{2.2em}
\renewcommand{\@tocrmarg}{3em}
\makeatother

% ---- 框 ----
\newenvironment{chapterintro}{\par\color{qmuted}}{\par\vspace{0.4em}}

% 題目（英文原題）。跨頁時續段右上角標「（續）」，免得看成新的一題（DE CH11 版面審查）。2026-10-27
\newtcolorbox{statement}{enhanced,breakable,frame hidden,arc=2pt,
  colback=qgraybg,left=10pt,right=10pt,top=6pt,bottom=6pt,
  before skip=0.2em,after skip=0.9em,before upper={\statementfont},
  overlay middle and last={\node[anchor=north east,inner sep=0pt,font=\scriptsize,text=qmuted]
    at ([xshift=-6pt,yshift=-4pt]frame.north east) {（續）};}}

% 結果框；預設無標籤，需要時 \begin{resultbox}[答案 (b)]
\newtcolorbox{resultbox}[1][]{enhanced,frame hidden,sharp corners,
  colback=qorangebg,borderline west={2pt}{0pt}{qorange},
  left=10pt,right=10pt,top=5pt,bottom=5pt,
  before skip=0.6em,after skip=0.8em,
  before upper={\if\relax\detokenize{#1}\relax\else{\small\bfseries\color{qorange}#1}\par\fi}}

% 提醒或補充
\newtcolorbox{notebox}[1][說明]{enhanced,breakable,frame hidden,sharp corners,
  colback=qblue!5,borderline west={2pt}{0pt}{qblue},
  left=11pt,right=11pt,top=6pt,bottom=6pt,
  before upper={{\small\bfseries\color{qblue}#1}\par}}

\tcbset{highlight math style={enhanced,frame hidden,sharp corners,colback=qorangebg,
  borderline west={2pt}{0pt}{qorange},left=6pt,right=6pt,top=3pt,bottom=3pt}}

% ---- 數學巨集（各專案共用的兩個；其餘記號在專案的 .sty）----
\newcommand{\dd}{\mathop{}\!\mathrm{d}}
\newcommand{\e}{\mathrm{e}}
% \keyresult{...} 的內容當文字排（數學用 $...$），放進結果框；已在數學環境內時改為橘色底
\newcommand{\keyresult}[1]{%
  \ifmmode\tcbhighmath{\displaystyle #1}%
  \else\begin{resultbox}{\centering #1\par}\end{resultbox}\fi}
